% 第 67 章骨架：导引框已定，正文待写。写作计划见 docs/PLAN.md“第 67 章”，写法见 docs/STYLE.md。
\chapter{安倍经济学与收益率曲线控制}\label{ch:abenomics}

\begin{chapterguide}
\guidequestion{日本央行怎样直接给长期利率设上限？退出为什么这么难？}
\begin{guidegoals}
\item 描述 2013 年起的量化质化宽松与“三支箭”；
\item 解释 2016 年收益率曲线控制与负利率的设计；
\item 评估它对通胀、汇率与财政的影响；
\item 追踪 2024 年退出收益率曲线控制以后的利率正常化，以及 2025—2026 年的财政争论。
\end{guidegoals}
\guideprereq{\chref{ch:zlb}\enspace\nameref{ch:zlb}；\chref{ch:japanbubble}\enspace\nameref{ch:japanbubble}}
\end{chapterguide}

\chapterstub
